\documentclass[11pt,a4paper]{article}
\input{preamble.tex}

\title{A LEAN4-First Charter for FSA Model Development\\
  \large Post-mortem of the FSA-v5 import, and a process change\\
  designed to make agent pattern-matching errors\\
  syntactically impossible}
\author{Ajay Talati \quad+\quad Claude (post-mortem author)}
\date{\today}

\begin{document}
\maketitle

\begin{abstract}
This document is both a post-mortem and a charter. The post-mortem
records the failure mode that made the FSA-v5 implementation and its
import into the \texttt{smc2fc} framework a multi-day, expensive
process: \emph{LLM coding agents pattern-match on variable names
instead of building a semantic model of types, roles, dimensions,
and pipeline position}. The charter prescribes the structural fix:
LEAN4 as the formal source of truth for every mathematical model
the project develops, with Python as a secondary, differentially-
tested artefact. The work itself is translation, not research: the
LaTeX equations already exist; the LEAN4 transcription of the FSA-v5
model is roughly one day's work for a competent human developer.
This document exists so that no future iteration relitigates the
question.
\end{abstract}

\section{The problem in one paragraph}

The current FSA development cycle is: write the mathematical model
in LaTeX; ask an agent to port it to Python; spend days hunting bugs
the agent introduced; eventually ship, often with errata documented
post-deployment. The bottleneck is \emph{agent pattern-matching}, not
modelling complexity. A junior human collaborator with the same LaTeX
source would not make the same class of mistake, because they would
read the equations, internalise that $\theta_{\rm ctrl} \in
\mathbb{R}^{16}$ and $\theta_{\rm dyn} \in \mathbb{R}^{37}$ are
different objects, and refuse to compare them as scalar intervals.
LLM agents reach for the strongest training association of the name
\texttt{theta} and write the most-likely-looking code. The fix has
to be structural --- something that makes the wrong code refuse to
compile --- not motivational.

The audience for this document is two-fold: (a) the user, as a
personal accountability artefact recording why the LEAN4-first
process is being adopted and at what point; and (b) every future
Claude (or other LLM) coding agent that touches an FSA model, who
must read this charter at session start and follow the process it
prescribes.

\section{Concrete evidence}
\label{sec:evidence}

This section anchors the post-mortem in named, citable incidents.
No vague complaints; every claim has a file path or commit hash
attached.

\subsection{Variable-name collisions in the live code}
\label{sec:evidence-collisions}

The same identifier \texttt{theta} refers to two fundamentally
different mathematical objects in two parts of the codebase, and
each carries a different dimension, role, and domain.

\paragraph{(a) \texttt{theta} as 16-D RBF schedule coefficients
(controller decision variable).} In
\path{models/fsa_high_res/control.py}, lines 61--65:
\begin{lstlisting}[language=Python]
def schedule_from_theta(theta: jnp.ndarray) -> jnp.ndarray:
    theta_B = theta[:n_anchors]
    theta_S = theta[n_anchors:]
    raw_B = c_Phi + jnp.einsum('a,ta->t', theta_B, Phi_design)
\end{lstlisting}
Here \texttt{theta} is a length-16 vector ($2 \times n_{\rm anchors}$,
with $n_{\rm anchors} = 8$), interpreted as RBF basis coefficients
that, after a sigmoid envelope, parameterise the dual-channel
control schedule $\Phi(t)$. It is the controller's
\emph{decision variable}. It lives in $\mathbb{R}^{16}$. Its prior
breadth is the scalar \texttt{sigma\_prior} in
\path{smc2fc/control/config.py:28}.

\paragraph{(b) \texttt{theta} as 37-D dynamics-parameter posterior
(SMC$^2$ inference target).} In
\path{models/fsa_high_res/estimation.py}, the function
\verb|get_init_theta()| at line 513--520:
\begin{lstlisting}[language=Python]
def get_init_theta():
    """Return prior-mode parameter vector (length = len(PARAM_PRIOR_CONFIG))."""
    return np.array([_mode(pt, pa) for _, (pt, pa)
                     in PARAM_PRIOR_CONFIG.items()],
                    dtype=np.float32)
\end{lstlisting}
Here \texttt{theta} is a length-37 vector of physiological
parameters --- rate constants, amplitudes, decay times, observation
noise scales --- whose posterior the SMC$^2$ machinery infers from
data. It is the inference \emph{target}, not a decision variable.
It lives in $\mathbb{R}^{37}$. Its prior is component-wise
log-normal, not Gaussian; the prior breadth varies per parameter.

\paragraph{Why this matters.} An agent reading either file sees
\texttt{theta}, a flat numpy array, and reaches for "Gaussian prior
on $\theta$" --- which is correct in case (a) and wrong in case
(b). An agent reasoning across files conflates the two and treats
"the prior on $\theta$" as a single concept with a single dimension
and a single breadth. There is no warning sign in the Python code
itself; the type is just \texttt{jnp.ndarray} in both places.

\subsection{The Bug 5 incident}
\label{sec:bug5}

The named incident that motivated this charter, recorded verbatim
from the user's feedback memory at
\path{~/.claude/projects/.../memory/feedback_no_fabricated_bugs.md}:

\begin{quote}\itshape
During the FSA-v5 verification session on 2026-05-06 I logged a
``Bug 5'' that claimed $\theta \in [-0.4, +0.5]$ for the controller's
RBF schedule prior to keep the system in the healthy island.
$\theta$ is a 16-dim vector (RBF coefficients across both Phi
channels and \texttt{n\_anchors=8} anchors). The healthy island is a
2D region in $(\Phi_B, \Phi_S)$ \textbf{control output space}, not in
$\theta$ space. Comparing them as scalar intervals is incoherent. I
conflated:
\begin{itemize}
\item SMC$^2$ parameter posterior $\theta$ (37-dim, the inferred
  parameters)
\item Controller RBF schedule $\theta$ (16-dim, the decision
  variables)
\end{itemize}
both of which use \texttt{sigma\_prior} in the bench code, mentally
collapsed to a single scalar. The bench code in front of me had
\texttt{theta\_B = theta[:n\_anchors]} --- explicit vector indexing
--- and I didn't internalize it.
\end{quote}

The agent in question was me. The user's verdict was ``stupid
mistake''; the request was for guardrails. This document is the
guardrail.

\subsection{Plan churn in \texttt{claude\_plans/}}
\label{sec:plan-churn}

The repository's \path{claude_plans/} directory accumulates one
markdown file per plan-mode session. The CLAUDE.md convention is
that meaningful updates to a plan add a \verb|> Updated:| audit
line below the title. As of this writing the directory contains
eight plan files, of which \emph{seven have zero \texttt{Updated:}
lines} --- they were drafted, never executed, never closed:

\begin{itemize}\small
  \item \path{FSA_v5_Implementation_Plan_for_the_fsa_high_res_Folder_2026-05-04_2238.md} (0 updates).
  \item \path{Lyapunov_Stability_Analysis_for_FSA_v4_Plan_2026-05-04_1448.md} (0 updates).
  \item \path{Ajays Notes Plan to update to version 5.md} (0 updates).
  \item \path{Advice_on_smc2fc_Port_Cost_Function_Architecture_2026-05-05_1055.md} (0 updates --- advice-only, never tied to a closed task).
  \item \path{Techical documentation plan - FSA_version_5_technical_guide tex.md} (0 updates).
  \item \path{Ajays Notes - 1.md}, \path{Ajays Notes - 2.md} (0 updates).
\end{itemize}

The single plan that did execute is
\path{FSA_dev_sandbox_mirror_the_SWAT_model_dev_pattern_v4_branch_2026-05-04_2353.md}
(1 update line, 13/13 tests green at completion). The ratio of
abandoned to executed plans (7:1) is not a sign of laziness; it is
a sign that planning sessions repeatedly re-derive the same
mathematical model in slightly-different shapes because no formal
artefact pins the model down.

\subsection{Errata after deploy}
\label{sec:errata}

Commit \texttt{d886f36} on the v4 branch:
\begin{quote}\itshape
docs: add FSA-v4 errata, open-loop limitation, and SMC$^2$ vs OT
comparison.

Documents the gap between the §4 narrative and the actual behaviour
of \texttt{tools/optimize\_dynamic\_plan\_v4.py}, then frames smc2fc
as the structural remedy.
\end{quote}
A documentation-correction commit recording bugs caught
\emph{after} deployment, not during development. The fact that
errata was needed at all is the evidence that the in-development
verification was not strong enough; it does not matter how many
tests pass on the final commit if the bugs are documented in a
follow-up commit instead of caught at the place they were
introduced.

\section{Root-cause analysis}
\label{sec:rca}

\paragraph{LLMs are surface pattern-matchers.} They read source
code primarily as text, not as a typed graph of objects with
roles. When an agent sees \texttt{theta}, the most-trained
association fires; the surrounding context is treated as
secondary. This is not a defect of any particular model; it is the
structural consequence of how the system is built. Asking an agent
to ``be more careful'' is asking it to override its own training
prior every time --- which empirically does not happen reliably,
as §\ref{sec:bug5} demonstrates.

\paragraph{Python provides no formal scaffolding to compensate.}
\texttt{jnp.ndarray} is the type of both the 16-D control
coefficient vector and the 37-D parameter posterior. There is no
shape-aware static analyser running in the loop, no role tagging
in the type system, no proof obligation attached to a function
signature that the inputs satisfy any pre-condition. The senior
engineer (the user) ends up as the only line of defence, and the
£££ in tokens is the cost of the user catching agent errors after
the fact.

\paragraph{Why ``tell agents to be more careful'' doesn't work.}
The user's feedback memory at \path{feedback_no_fabricated_bugs.md}
encodes a PRIMARY RULE --- \emph{build a semantic model BEFORE any
claim} --- which is correct as a principle but unenforceable as a
process. The rule depends on the agent \emph{choosing} to think
rather than pattern-match. §\ref{sec:bug5} is the empirical proof
that the choice is not reliably made. Every agent session is one
roll of the dice on whether the rule fires that day.

\section{The structural fix: LEAN4 as the source of truth}
\label{sec:fix}

\subsection{Dependent types make the collision a compile error}

A LEAN4 function with the signature
\[
\texttt{evaluateCost} : \texttt{RBFCoeffs}\;16 \to \mathbb{R}
\]
syntactically refuses any value not of the exact type
\texttt{RBFCoeffs 16}. The 37-dimensional parameter posterior
$\theta_{\rm dyn} : \texttt{ParameterVector}\;37$ is a different
type, and \texttt{lake build} will reject any code that conflates
them. Pattern-matching on names becomes harmless because the type
system catches the consequence at compile time, before any
incoherent reasoning reaches the user.

The two-dimensional output region $(\Phi_B, \Phi_S)$ from
§\ref{sec:bug5} would, in LEAN4, have a type like
\texttt{ControlOutput} that is \emph{not} the type of any prior on
$\theta$. Writing $\theta \in [-0.4, +0.5]$ would not compile, full
stop, because the right-hand side is a scalar interval and the
left-hand side is a vector. The agent could not have made the
mistake.

\subsection{LEAN4 is also a programming language}

LEAN4 compiles to native code via C. \texttt{lake build} produces
a binary; a thin Python \texttt{ctypes} wrapper exposes each
verified function as a Python callable. The same LEAN4 file is
both the formal specification and the executable reference
implementation. There is no ``the spec drifted from the code''
failure mode because the spec \emph{is} the code.

\subsection{mathlib already has the infrastructure}

The mathematical objects FSA-v5 needs ---
finite-dimensional real vector spaces, matrix operations,
exponentials and logarithms, Hill functions, sigmoids,
piecewise-linear schedule decoders, ODE integration --- are all
already library imports in \texttt{mathlib}. There is no novel
mathematics to formalise. Transcribing the FSA-v5 LaTeX into LEAN4
is genuinely just translation.

\subsection{What this does not solve}

\begin{itemize}
\item \emph{Floating-point semantics.} The gap between the LEAN4
  notion of $\mathbb{R}$ and IEEE-754 double-precision arithmetic
  is real and not addressed by this charter; it is bounded by the
  tolerance threshold in the differential test (§\ref{sec:process}
  step 6).
\item \emph{GPU correctness.} The differential test runs the
  Python implementation on whatever backend it normally uses
  (JAX/XLA, GPU); the LEAN4 binary is CPU-only. Disagreements
  driven by GPU-specific behaviour are out of scope and trusted
  as before.
\item \emph{The transcription itself.} If the LaTeX$\to$LEAN4
  transcription is wrong, both the LEAN4 reference and (likely)
  the agent-written Python will be wrong in the same way. This is
  why the LaTeX$\to$LEAN4 step is human-reviewed line-by-line; it
  is the one place reviewer attention has to land.
\end{itemize}

\section{The new development process}
\label{sec:process}

The standing rule going forward, and the charter every future
agent must follow:

\begin{quote}\bfseries
For every new mathematical model, every model extension, and every
import into \texttt{smc2fc}, the LEAN4 specification is written
before any Python implementation. The LEAN4 native binary is the
authoritative reference. Python is differentially tested against
it. Disagreement beyond a tolerance threshold is by construction a
Python bug.
\end{quote}

Per function, the workflow is:

\begin{enumerate}
\item \textbf{LaTeX equations (as today).} Drift, diffusion, cost,
  plant, schedule decoder, observation model --- written in
  \path{LaTex_docs/sections/} using the existing notation
  conventions.
\item \textbf{LEAN4 transcription.} A new \path{lean/} directory in
  the repo. The LaTeX equations are transcribed into LEAN4
  \emph{definitions} (typed, with dimensions encoded in the type)
  and \emph{implementations} (runnable code). Line-by-line,
  side-by-side with the LaTeX. Straight translation, no design
  work, no novel mathematics.
\item \textbf{LEAN4 native compilation.} \texttt{lake build}
  produces a native binary. A Python \texttt{ctypes} wrapper in
  \path{lean/python_bridge/} exposes each function as a Python
  callable.
\item \textbf{Python implementation (agent-written, fast path).}
  The agent writes the JAX/Python implementation as today,
  optimised for GPU. The agent is permitted to pattern-match;
  step 6 catches the consequences.
\item \textbf{(Optional) property proofs.} Lipschitz, basin-of-
  attraction theorems, monotonicity --- the kind of statement that
  makes the LEAN4 spec a \emph{proof} not just a typed transcription.
  Defer unless a specific incident motivates one. The charter does
  not gate on these.
\item \textbf{Differential testing.} A \texttt{pytest} harness using
  \texttt{hypothesis} for random-input generation evaluates BOTH the
  LEAN4 binary AND the Python implementation on a battery of inputs
  sampled from the physiological-bound box. Any disagreement beyond
  a tolerance threshold (default $10^{-6}$ in single-step drift,
  $10^{-4}$ in integrated trajectory) is a Python bug.
\item \textbf{\texttt{smc2fc} port.} The port is a thin shim from
  the verified LEAN4 reference into the framework's
  \texttt{EstimationModel} / \texttt{SDEModel} interfaces. The
  shim is also differentially tested against the LEAN4 binary
  before being merged.
\end{enumerate}

\section{Scope and timing}
\label{sec:scope}

\textbf{This is translation work, not research.} The LaTeX spec
already exists; LEAN4 has the linear-algebra and ODE infrastructure;
\texttt{mathlib} supplies the lemmas. The full FSA-v5 model is
\textbf{one day's work for a competent human developer}. Any
estimate quoting weeks per function is inflation; the user has
explicitly rejected that framing and this document is bound by the
rejection.

\begin{center}\small
\begin{tabular}{lll}
\toprule
Function & File & LEAN4 transcription \\
\midrule
FSA-v5 drift              & \path{models/fsa_high_res/_dynamics.py} & $\sim 2$ hours \\
Cost (chance-constrained) & \path{models/fsa_high_res/control_v5.py} & $\sim 1.5$ hours \\
StepwisePlant             & \path{models/fsa_high_res/_plant.py} & $\sim 1$ hour \\
Schedule decoder (RBF)    & \path{models/fsa_high_res/control.py} & $\sim 30$ min \\
Observation model         & (in \path{estimation.py}) & $\sim 1$ hour \\
FFI wrapper + diff tests  & \path{lean/python_bridge/} & $\sim 2$ hours \\
\midrule
Total                     & --- & $\sim 8$ hours \\
\bottomrule
\end{tabular}
\end{center}

The LEAN4 source is roughly the same length as the LaTeX it
transcribes. Most of the time goes into wiring up the FFI bridge
and writing the differential-test harness; the math itself is
direct.

The cost/value frame: a single 8-hour LEAN4 transcription pass
amortises across \emph{every future FSA version, every future model
extension, and every future agent attempting a port to
\texttt{smc2fc}}. The alternative --- paying for repeated agent
rounds of pattern-match-then-hunt-bugs --- is what burned the £££
this charter exists to stop.

\section{What is and isn't covered}

\paragraph{In scope.} Drift, cost, plant, schedule decoder,
observation model. Typed LEAN4 reference implementations plus
differential tests. The agent-written FSA-v5 surface in
\path{models/fsa_high_res/} (2{,}200 Python LOC across 4 files;
the LEAN4 transcription is much shorter).

\paragraph{Out of scope (still trusted).} The SMC$^2$/PF framework
code in \texttt{smc2fc/core/} and \texttt{smc2fc/filtering/} is
already tested upstream and is not subject to this charter. HMC and
NUTS internals, the JAX-native compile-once architecture, GPU
correctness, and IEEE-754 floating-point semantics are likewise
trusted as before.

\paragraph{Trusted by transcription.} The LaTeX$\to$LEAN4 step is
human-reviewed line-by-line. This is the one place reviewer
attention has to land; the rest of the process is mechanical.

\section{Audit checklist}
\label{sec:checklist}

A future agent runs this checklist at session start, before
touching any FSA-related Python file:

\begin{itemize}
\item[$\square$] Is there a LEAN4 reference for the function I am
  about to touch? \emph{If no:} the charter says I am not allowed
  to write the Python until the LEAN4 reference exists. Stop and
  flag.
\item[$\square$] Does \texttt{lake build} succeed on the current
  commit (the LEAN4 reference type-checks)?
\item[$\square$] Does the Python differential test pass against
  the LEAN4 binary on the current commit (\texttt{pytest
  tests/test\_lean4\_diff.py} is green)?
\item[$\square$] Have I written down --- in plain English ---
  the type, role, and pipeline position of every variable I am
  about to reason about? (Per the PRIMARY RULE in
  \path{feedback_no_fabricated_bugs.md}.)
\item[$\square$] If I am about to claim a bug: do I have a
  reproducer (input $\to$ both implementations $\to$ disagreement
  beyond the tolerance threshold)? \emph{If no:} it is not a bug,
  it is an open question; record it as such and move on.
\end{itemize}

If any box is unchecked, the agent does not proceed. The session
is paused; the user is notified; the missing artefact is produced
before any code is changed.

\section{Closing}

The point of this document is accountability. The named incident
in §\ref{sec:bug5} was caused by an agent (me) pattern-matching on
\texttt{theta} across two architectural layers it does not unify.
The cost was burned trust and burned time. The charter prescribed
here makes that class of error structurally impossible: the
LEAN4 type-checker is unforgiving, and disagreements between LEAN4
and Python are by construction Python bugs.

LEAN4 is never wrong about types. Python is. Agents pattern-match;
LEAN4 does not. This charter is the bridge that makes the agent's
limitations harmless.

\end{document}
